\documentclass{article}
\usepackage[margin=1in]{geometry}
\usepackage{xcolor}
\usepackage{amsmath}
\usepackage{amssymb}
\usepackage{fontspec}
\usepackage{tikz}
\usetikzlibrary{tikzmark}
\usepackage{enumitem}

\definecolor{bgdark}{HTML}{121212}
\definecolor{textlight}{HTML}{E0E0E0}
\definecolor{subtext}{HTML}{A0A0A0}

\pagecolor{bgdark}
\color{textlight}
\pagestyle{empty}

% #1 = number, #2 = node name used later for the connector lines
\newcommand{\stepnum}[2]{%
  \tikzmarknode{#2}{%
    \begin{tikzpicture}[baseline=(char.base)]
      \node[draw=textlight, circle, inner sep=1.5pt, font=\sffamily\small] (char) {#1};
    \end{tikzpicture}%
  }%
}

\begin{document}
\sffamily

\begin{description}[leftmargin=2em]

  \item[\stepnum{1}{s1}] \textbf{Decompose each number into tens and ones}\\[0.3em]
  {\color{subtext}\textit{Separating values makes mental math simple}}
  \begin{itemize}
    \item $11  = 10  + 1$
    \item $51  = 50  + 1$
    \item $101 = 100 + 1$
  \end{itemize}
  \vspace{1em}

  \item[\stepnum{2}{s2}] \textbf{Group and add all the tens together}\\[0.8em]
  Combine the multiples of 10:
  \[
    10 + 50 + 100 = 160
  \]
  \vspace{1em}

  \item[\stepnum{3}{s3}] \textbf{Group and add all the remaining ones}\\[0.8em]
  Combine the ones:
  \[
    \mathbf{1} + \mathbf{1} + \mathbf{1} = \mathbf{3}
  \]
  \vspace{1em}

  \item[\stepnum{4}{s4}] \textbf{Add the totals together}\\[0.8em]
  Combine the sum of the tens with the sum of the ones:
  \[
    160 + \mathbf{3} = \mathbf{163}
  \]

\end{description}

\begin{tikzpicture}[remember picture, overlay]
  \draw[dotted, thick, textlight!60] (s1.south) -- (s2.north);
  \draw[dotted, thick, textlight!60] (s2.south) -- (s3.north);
  \draw[dotted, thick, textlight!60] (s3.south) -- (s4.north);
\end{tikzpicture}

\end{document}